\section{Methodology: Principled Communication via Active Inference}
\label{sec:methodology}

We propose a novel framework for multi-agent coordination grounded in the Free Energy Principle (FEP) \cite{friston2010free, friston2017active}. Unlike traditional Multi-Agent Reinforcement Learning (MARL) that treats communication as an auxiliary signal to maximize a shared reward $\mathcal{R}$ \cite{lowe2017multi, rashid2018qmix}, our framework, \textbf{ActInf-Comm}, casts communication as a fundamental epistemic action designed to minimize Variational Free Energy (VFE). This shift allows agents to resolve non-stationarity by actively inferring the latent strategies of competitors, a critical requirement in high-frequency environments like Real-Time Bidding (RTB).

\begin{figure*}[t]
    \centering
    % \includegraphics[width=\textwidth]{fig1_framework.pdf}
    \caption{\textbf{The ActInf-Comm Framework Architecture.} (Left) The hierarchical generative model where Level 0 ($L_0$) processes immediate auction dynamics and Level 1 ($L_1$) infers market regimes. (Center) The Action Selection loop where policies $\pi$ are evaluated by Expected Free Energy ($\mathcal{G}$), decomposing into Epistemic Value (information gain) and Instrumental Value (preference satisfaction). (Right) The LLM-augmented inference engine that maps high-dimensional market observations into semantic tokens for accelerated belief updating. Communication ($\alpha_c$) serves to synchronize these latent beliefs across agents, preventing collusive traps by prioritizing model transparency over reward exploitation.}
    \label{fig:framework}
\end{figure*}

\subsection{The Active Inference Agent: Generative Model and VFE}
\label{sec:agent_model}

Each agent $i \in \{1, \dots, N\}$ is modeled as a Bayesian inference engine that maintains a generative model $P(o, s, \pi)$ over observations $o$, hidden states $s$, and policies $\pi$. The hidden states $s = \{s_{env}, s_{other}\}$ encompass both the environmental dynamics (e.g., market price distributions) and the latent intentions of other agents (e.g., budget constraints, risk profiles).

The agent's objective is to minimize the Variational Free Energy $\mathcal{F}$, which serves as an upper bound on the surprise (negative log-evidence):
\begin{equation}
    \mathcal{F}(q) = \underbrace{\mathbb{D}_{KL}[q(s) \parallel p(s)]}_{\text{Complexity}} - \underbrace{\mathbb{E}_{q(s)}[\ln p(o|s)]}_{\text{Accuracy}} \approx -\ln p(o)
\label{eq:vfe}
\end{equation}
where $q(s)$ is the agent's internal belief (recognition density) about the world. By minimizing $\mathcal{F}$, the agent simultaneously optimizes its internal representation of the non-stationary environment and its predictions of future observations \cite{parr2022active}. 

Policy selection in ActInf-Comm is governed by the \textit{Expected Free Energy} (EFE), denoted as $\mathcal{G}(\pi)$. For a future time horizon $\tau$, the agent evaluates policies $\pi$ based on:
\begin{equation}
    \mathcal{G}(\pi) = \sum_{\tau} \mathbb{E}_{q(o_\tau, s_\tau | \pi)} [\ln q(s_\tau | \pi) - \ln p(o_\tau, s_\tau)]
\label{eq:efe}
\end{equation}
This objective naturally decomposes into \textit{instrumental value} (seeking preferred outcomes, e.g., winning an auction at low cost) and \textit{epistemic value} (seeking information to reduce uncertainty about $s$). In non-stationary MAS, the epistemic component is vital, as it compels agents to explore and update their models of competitors' shifting strategies before committing to intensive resource allocation \cite{da2020active}.

\subsection{Communication as an Epistemic Action}
\label{sec:comm_action}

In our framework, the action space $\mathcal{A}$ is partitioned into environmental actions $\mathcal{A}_e$ (e.g., bid prices in RTB) and communicative actions $\mathcal{A}_c$ (e.g., broadcasted signals). Crucially, $\mathcal{A}_c$ is not rewarded by the environment; instead, its utility emerges from its ability to reduce the expected surprise of the collective system.

We define a communicative action $\alpha_c \in \mathcal{A}_c$ as an ``epistemic probe.'' When agent $i$ emits a signal, it intends to modify the belief $q_j(s)$ of agent $j$, thereby making agent $j$'s future actions more predictable. The epistemic value of communication is formulated as the mutual information between the signal and the latent states:
\begin{equation}
    \text{EpistemicValue}(\alpha_c) = \mathbb{E}_{q(o_\tau | \alpha_c)} [ \mathbb{D}_{KL}[q(s_\tau | o_\tau, \alpha_c) \parallel q(s_\tau | \alpha_c)] ]
\end{equation}
Unlike standard MARL signals that often converge to ``cheap talk'' or collusive price-fixing in RTB \cite{zheng2022multiagent, brown2020language}, ActInf agents communicate to maintain ``generalized synchrony'' \cite{friston2015active}. Because the EFE prioritizes model accuracy, agents are disincentivized from sending deceptive signals that would increase their own long-term VFE by making the environment's response to their collusion unpredictable. This provides a principled, information-theoretic barrier against emergent collusion.

\subsection{Hierarchical Generative Models for Non-Stationary MAS}
\label{sec:hierarchical_model}

To handle the multi-scale temporal dynamics of RTB, we employ a Hierarchical Active Inference (HAI) structure. The model is organized into two primary levels:

\begin{enumerate}
    \item \textbf{Level 0 ($L_0$): Fast-Scale Dynamics.} This level operates at the frequency of individual auctions. It performs inference on immediate states $s^{(0)}$, such as current bid requests and winning price volatility. The transition dynamics $\mathcal{P}^{(0)}$ are conditioned on the higher-level state $s^{(1)}$.
    \item \textbf{Level 1 ($L_1$): Slow-Scale Regimes.} This level infers the ``context'' or ``regime'' of the market (e.g., high-competition periods, end-of-quarter budget flushing). Changes at this level represent the source of non-stationarity for $L_0$.
\end{enumerate}

The interaction between levels is governed by the exchange of prediction errors ($\epsilon = o - \hat{o}$) flowing upward and predictions flowing downward. In the context of RTB, an unexpected jump in market prices generates a large prediction error at $L_0$, which forces $L_1$ to update its belief about the current market regime. Communication actions $\alpha_c$ are often selected at $L_1$ to synchronize these regime beliefs across agents, ensuring coordination without the need for exhaustive re-training \cite{millidge2023mathematical}.

\subsection{LLM-Augmented Inference and Emergent Concept Representation}
\label{sec:llm_inference}

A significant challenge in applying ActInf to complex MAS is the computational complexity of evaluating $\mathcal{G}(\pi)$ over high-dimensional state spaces. We address this by integrating Large Language Models (LLMs) as \textit{amortized inference engines} \cite{hu2023language, wang2024active}. 

We utilize a pre-trained LLM (e.g., GPT-4o) to map sequences of market observations and agent signals into a compact semantic latent space $\mathcal{Z}$. The LLM serves two purposes:
\begin{itemize}
    \item \textbf{State Abstraction}: It translates raw auction logs into high-level semantic descriptions (e.g., ``competitor A is exhibiting predatory pricing behavior''), which are then used as observations for the $L_1$ generative model.
    \item \textbf{Concept Synthesis}: When agents encounter novel non-stationary regimes not seen during training, the LLM-augmented component can propose new semantic tokens to represent these states. This allows ActInf agents to communicate \textit{unencoded concepts}—abstract strategies that were not part of the initial state space $\mathcal{S}$—by leveraging the world knowledge embedded in the LLM \cite{wei2022chain}.
\end{itemize}

Specifically, the LLM generates a distribution over possible latent explanations $q(s | \text{history}, \text{context})$, which acts as a sophisticated prior for the VFE minimization process. This acceleration allows the agents to maintain real-time responsiveness in milliseconds-level RTB auctions while performing deep hierarchical inference. By grounding the LLM's outputs in the rigorous VFE objective, we ensure that the ``hallucinations'' typical of generative models are filtered out by the accuracy requirements of the generative model $\mathcal{F}$ \cite{fountas2020active}.